\documentclass[12pt]{article}

\usepackage[utf8]{inputenc}
\usepackage[margin=1in]{geometry}
\usepackage{setspace}
\usepackage{url}

\doublespacing
\setlength{\emergencystretch}{3em}

\title{Learning the Shape of a Responsible Path:\\
An Identity Memoir on Alice Munro's ``Boys and Girls''}

\author{%
  Ma Toan Bach\\
  Student Number: 112708227\\
  LSO100KCC\\
  Short Story \& Identity Memoir
}

\date{}

\begin{document}

\maketitle

When I think about identity, I do not think first about one dramatic decision. I think about ordinary conversations that slowly teach a person what kind of future is considered acceptable. For me, one important part of growing up was learning how to respond to expectations about school, work, and responsibility. I was often encouraged to choose a path that looked stable, practical, and safe. I understood why people cared about this. They wanted me to have security and not waste opportunities. Still, I also felt the pressure of being measured by a future that other people had already imagined for me. Alice Munro's ``Boys and Girls'' helps me understand that identity is not only something we discover privately. It is also shaped by family, social roles, and the quiet pressure to become the person others expect.

In my teenage years, the question of the future started to feel less like an open question and more like a test. I remember one afternoon when I sat with a course-selection form in front of me. The page looked simple: boxes, program names, deadlines, and blank spaces waiting for my choices. Around me, the advice sounded reasonable. Choose something practical. Think about job security. Do not waste time on a path that might not lead anywhere clear. I held the pen and tried to look confident, but I felt as if the form was asking more than what courses I wanted. It was asking what kind of person I was allowed to become.

The pressure was not harsh, which made it harder to name. No one needed to force me directly, because concern can be persuasive. Advice came from people who cared about me and wanted me to succeed. I listened because I knew they were not trying to hurt me. I nodded, said that I understood, and kept my doubts mostly to myself. Inside, however, I was beginning to feel a distance between the life being described for me and the person I was becoming.

Munro's narrator feels a similar tension, even though her situation is different from mine. At the beginning of ``Boys and Girls,'' she takes pride in helping her father with the foxes. She carries water, works outside, and feels useful in a world that seems important. When her father introduces her as his ``new hired hand,'' she turns away ``red in the face with pleasure'' (Munro, 1968, p.~5). That moment matters because it shows how identity can grow from recognition. The narrator does not only want to help; she wants her work to mean something. She wants to be seen as capable.

I recognize that feeling. When I did well in school or completed work that mattered to me, I wanted the same kind of recognition. I wanted to feel that my effort proved something about who I was. But, like Munro's narrator, I also learned that recognition can be temporary if it does not fit the role others expect. In the story, the narrator's mother says that when Laird is older, he will be ``a real help,'' and then says she can use the narrator more in the house (Munro, 1968, p.~6). The narrator realizes that her place beside her father may not be permanent. Her work matters, but gender changes how other people understand that work.

My own experience was not about being told to become a girl in the way Munro's narrator is. The expectation placed on me was different: to become practical, predictable, and careful. Still, the structure of the pressure felt familiar. Other people were not only advising me about choices; they were also describing the kind of person I should become. I began to understand that identity can be shaped through repetition. If people repeatedly describe one path as sensible and another as unrealistic, even personal desire starts to sound irresponsible.

Munro makes this process especially clear when the narrator begins to understand the word ``girl'' differently. She says that it once seemed ``innocent and unburdened,'' but later becomes ``a definition, always touched with emphasis, with reproach and disappointment'' (Munro, 1968, p.~8). I think this is one of the most powerful insights in the story. A word that should simply describe someone becomes a limit. For me, words like ``safe,'' ``stable,'' and ``responsible'' sometimes worked in the same way. They were good words, but they could also become boundaries. They suggested that choosing differently was not only risky but immature.

The central memory I connect to this story is not a dramatic argument. It is a quieter memory of sitting with my future reduced to options on a form and suggestions from other people. I remember the pressure of wanting to make everyone feel confident in me. I did not want to seem ungrateful or careless. I also did not want to disappoint people who had invested hope in my success. Yet I kept returning to the same private question: if I only choose the path that makes others comfortable, will I still recognize myself in that future?

That question did not produce an immediate answer. At first, I tried to solve the conflict by separating the two parts of myself. Outwardly, I acted practical. Inwardly, I protected the interests and questions that made me feel more like myself. I told myself that I could be responsible now and think about personal goals later. But later kept moving farther away. The more I postponed my own judgment, the easier it became to believe that my uncertainty was a weakness rather than a sign that I needed to think honestly.

This is why the Flora scene in Munro's story stays with me. When the narrator is told to shut the gate, she instead opens it wide and lets Flora run through. She says, ``I did not make any decision to do this, it was just what I did'' (Munro, 1968, p.~13). Her action does not save Flora in a practical sense, but it reveals something true about her. For a moment, she acts from sympathy, resistance, and identification with freedom. Later, she admits, ``I was on Flora's side,'' even though that makes her ``no use'' to her father's work (Munro, 1968, p.~14).

My own act of resistance was much smaller. It was not a single heroic choice. It was the slow decision to take my own uncertainty seriously. I began to ask what kind of work, study, and life would allow me to grow rather than only appear safe. I still respected practical advice, but I stopped treating it as the only language of maturity. That change was quiet, but it mattered. Like the narrator's action at the gate, it represented a private shift before it became something I could fully explain.

Writing about this now helps me see that growing up is not simply choosing between obedience and rebellion. The people who wanted me to choose responsibly were not enemies. In many ways, they were right to care about stability. Munro's story also complicates easy judgment. The painful part is that love and limitation can exist together. People can care about us and still imagine us too narrowly.

The ending of ``Boys and Girls'' shows the danger of accepting those narrow definitions. After the father says, ``She's only a girl,'' the narrator thinks, ``Maybe it was true'' (Munro, 1968, p.~16). That final sentence is painful because it shows how an outside judgment can become an inside belief. My experience taught me to be careful with that process. When I heard that one path was responsible and another was unrealistic, I sometimes began to believe that my own questions were childish. Munro's story helped me name that pressure and understand how quietly it works.

I do not think identity means ignoring everyone else's advice. I also do not think freedom means escaping all expectations. Instead, I have learned that identity requires reflection: deciding which expectations contain wisdom and which ones make the self smaller. A responsible path should not erase personal growth. It should include the courage to ask whether the future being offered actually belongs to you. Munro's narrator does not fully escape the definition placed on her, but her moment at the gate remains important because it shows that some part of her recognizes freedom. My own growth has been similar. I am still learning how to balance security with self-knowledge, but I understand now that becoming responsible should not mean becoming unrecognizable to myself.

\clearpage

\section*{References}

\hangindent=0.5in\noindent Munro, A. (1968). \textit{Boys and girls} [PDF]. \url{https://jerrywbrown.com/wp-content/uploads/2014/02/Boys-and-Girls-by-Alice-Munro.pdf}

\end{document}
